\begin{env}[4.1.6]
\label{0.4.1.6-fr}
Supongamos que $\mathcal{A}$ es un haz de anillos no necesariamente
conmutativos, que $\mathcal{J}$ es un haz de ideales a la izquierda de
$\mathcal{A}$ y que $\mathcal{F}$ es un $\mathcal{A}$-Módulo a la izquierda.
Se denota entonces por $\mathcal{J}\mathcal{F}$ el sub-$\mathcal{A}$-Módulo de
$\mathcal{F}$ que es la imagen de $\mathcal{J}\otimes_{\mathbf Z}\mathcal{F}$
(donde $\mathbf Z$ es el haz asociado al prehaz constante $U\to\mathbf Z$) por
la aplicación canónica $\mathcal{J}\otimes_{\mathbf Z}\mathcal{F}\to\mathcal{F}$.
Es claro que, para todo $x\in X$, se tiene
$(\mathcal{J}\mathcal{F})_x=\mathcal{J}_x\mathcal{F}_x$. Cuando
$\mathcal{A}$ es conmutativo, $\mathcal{J}\mathcal{F}$ es también la imagen
canónica de $\mathcal{J}\otimes_{\mathcal{A}}\mathcal{F}\to\mathcal{F}$.
Es inmediato que $\mathcal{J}\mathcal{F}$ es también el $\mathcal{A}$-Módulo
asociado al prehaz $U\to\Gamma(U,\mathcal{J})\Gamma(U,\mathcal{F})$. Si
$\mathcal{J}_1$, $\mathcal{J}_2$ son dos haces de ideales a la izquierda de
$\mathcal{A}$, se tiene
$\mathcal{J}_1(\mathcal{J}_2\mathcal{F})=(\mathcal{J}_1\mathcal{J}_2)\mathcal{F}$.
\end{env}
